% GENERATED FILE. Edit canonical lessons, then run npm run generate:books.
% canonical-source-hash: fnv1a64:16bc24b1311578a4
% canonical-lessons: UR-C103-hawaijahaz, UR-C103-relgari, UR-C103-pas, UR-C103-sakna, UR-C103-pana

\chapter{Aeroplane, Train, Near, To be able to, To manage to}
\label{ch:ur-aeroplane-train-near-to-be-able-to-to-manage-to}
\hlchaptermodality{103}

\begin{chapteropening}
\textbf{By the end of this chapter:} \emph{I can say an aeroplane, a train, near, to be able to and to manage to.}

Complete the last lesson of chapter 103: I can say an aeroplane, a train, near, to be able to and to manage to.
\end{chapteropening}

\section[\texorpdfstring{\ur{ہوائی} \ur{جہاز} (hawāī jahāz)}{ہوائی جہاز (hawāī jahāz)}]{\ur{ہوائی} \ur{جہاز} (hawāī jahāz) — an aeroplane}

\label{lesson:UR-C103-hawaijahaz}



\begin{quote}
Before the new one: say the Urdu for an office, then the Urdu for a library.
\end{quote}

\subsection*{You'll want to know: \ur{ہوائی} \ur{جہاز}}
\textbf{\ur{ہوائی} \ur{جہاز}} — \emph{hawāī jahāz} — \textquotedblleft{}an aeroplane\textquotedblright{}.

\begin{cousinweb}
\textbf{\ur{ہوائی} \ur{جہاز}} (\emph{hawāī jahāz}) is \textquotedblleft{}air-ship\textquotedblright{}: \emph{hawāī}, \textquotedblleft{}of the air\textquotedblright{}, from Arabic \emph{hawāʾ}, \textquotedblleft{}air\textquotedblright{}, and \emph{jahāz}, \textquotedblleft{}ship\textquotedblright{}, from Arabic \emph{jahāz}, \textquotedblleft{}equipment\textquotedblright{}, both through Persian (Lane, roots \emph{h-w-y}, \emph{j-h-z}).
\end{cousinweb}

\begin{grammarlens}[title={What you've built}]
One more word for this chapter.
\end{grammarlens}

\subsection*{Your turn}
\begin{itemize}
\raggedright
  \item \emph{Say it:} \emph{hawāī jahāz}
  \item \emph{Say it:} \emph{hawāī jahāz}, once more
  \item \emph{Say it:} say \emph{hawāī jahāz}
  \item \emph{Recall:} say the Urdu for an office, then the Urdu for a library, then say \emph{hawāī jahāz} again
\end{itemize}

\begin{tcolorbox}[breakable,colback=teal!4,colframe=teal!35!black,title={Before you move on}]
What does \ur{ہوائی} \ur{جہاز} mean? (\textquotedblleft{}An aeroplane\textquotedblright{}.) Say it once more.
\end{tcolorbox}
\section[\texorpdfstring{\ur{ریل} \ur{گاڑی} (rel gāṛī)}{ریل گاڑی (rel gāṛī)}]{\ur{ریل} \ur{گاڑی} (rel gāṛī) — a train}

\label{lesson:UR-C103-relgari}



\begin{quote}
Before the new one: say the Urdu for a library, then the Urdu for an aeroplane.
\end{quote}

\subsection*{You'll want to know: \ur{ریل} \ur{گاڑی}}
\textbf{\ur{ریل} \ur{گاڑی}} — \emph{rel gāṛī} — \textquotedblleft{}a train\textquotedblright{}.

\begin{cousinweb}
\textbf{\ur{ریل} \ur{گاڑی}} (\emph{rel gāṛī}) is \textquotedblleft{}rail cart\textquotedblright{}: \emph{rel} is a loan from English \emph{rail}, and \emph{gāṛī} goes back to a reconstructed \emph{*gāḍḍa}, \textquotedblleft{}cart\textquotedblright{}, which CDIAL links, tentatively, with Sanskrit \emph{garta}, \textquotedblleft{}seat of a chariot\textquotedblright{} (CDIAL 4116).
\end{cousinweb}

\begin{grammarlens}[title={What you've built}]
One more word for this chapter.
\end{grammarlens}

\subsection*{Your turn}
\begin{itemize}
\raggedright
  \item \emph{Say it:} \emph{rel gāṛī}
  \item \emph{Say it:} \emph{rel gāṛī}, once more
  \item \emph{Say it:} say \emph{rel gāṛī}
  \item \emph{Recall:} say the Urdu for a library, then the Urdu for an aeroplane, then say \emph{rel gāṛī} again
\end{itemize}

\begin{tcolorbox}[breakable,colback=teal!4,colframe=teal!35!black,title={Before you move on}]
What does \ur{ریل} \ur{گاڑی} mean? (\textquotedblleft{}A train\textquotedblright{}.) Say it once more.
\end{tcolorbox}
\section[\texorpdfstring{\ur{پاس} (pās)}{پاس (pās)}]{\ur{پاس} (pās) — near; with (someone)}

\label{lesson:UR-C103-pas}



\begin{quote}
Before the new one: say the Urdu for an aeroplane, then the Urdu for a train.
\end{quote}

\subsection*{You'll want to know: \ur{پاس}}
\textbf{\ur{پاس}} — \emph{pās} — \textquotedblleft{}near; with (someone)\textquotedblright{}.

\begin{cousinweb}
\textbf{\ur{پاس}} (\emph{pās}) is inherited. It comes from Sanskrit \emph{pārśve}, \textquotedblleft{}near\textquotedblright{}, a form of \emph{pārśva}, \textquotedblleft{}side\textquotedblright{} (CDIAL 8118).
\end{cousinweb}

\begin{grammarlens}[title={What you've built}]
One more word for this chapter.
\end{grammarlens}

\subsection*{Your turn}
\begin{itemize}
\raggedright
  \item \emph{Say it:} \emph{pās}
  \item \emph{Say it:} \emph{pās}, once more
  \item \emph{Say it:} say \emph{pās}
  \item \emph{Recall:} say the Urdu for an aeroplane, then the Urdu for a train, then say \emph{pās} again
\end{itemize}

\begin{tcolorbox}[breakable,colback=teal!4,colframe=teal!35!black,title={Before you move on}]
What does \ur{پاس} mean? (\textquotedblleft{}Near; with (someone)\textquotedblright{}.) Say it once more.
\end{tcolorbox}
\section[\texorpdfstring{\ur{سکنا} (saknā)}{سکنا (saknā)}]{\ur{سکنا} (saknā) — to be able to, can}

\label{lesson:UR-C103-sakna}



\begin{quote}
Before the new one: say the Urdu for a train, then the Urdu for near; with.
\end{quote}

\subsection*{You'll want to know: \ur{سکنا}}
\textbf{\ur{سکنا}} — \emph{saknā} — \textquotedblleft{}to be able to, can\textquotedblright{}.

\begin{cousinweb}
\textbf{\ur{سکنا}} (\emph{saknā}) is a tadbhava, a word inherited from Sanskrit and reshaped by sound change. It comes from \emph{śaknoti}, \textquotedblleft{}is able\textquotedblright{} (CDIAL 12252).
\end{cousinweb}

\begin{grammarlens}[title={What you've built}]
One more word for this chapter.
\end{grammarlens}

\subsection*{Your turn}
\begin{itemize}
\raggedright
  \item \emph{Say it:} \emph{saknā}
  \item \emph{Say it:} \emph{saknā}, once more
  \item \emph{Say it:} say \emph{saknā}
  \item \emph{Recall:} say the Urdu for a train, then the Urdu for near; with, then say \emph{saknā} again
\end{itemize}

\begin{tcolorbox}[breakable,colback=teal!4,colframe=teal!35!black,title={Before you move on}]
What does \ur{سکنا} mean? (\textquotedblleft{}To be able to, can\textquotedblright{}.) Say it once more.
\end{tcolorbox}
\section[\texorpdfstring{\ur{پانا} (pānā)}{پانا (pānā)}]{\ur{پانا} (pānā) — to manage to, to get}

\label{lesson:UR-C103-pana}



\begin{quote}
Before the new one: say the Urdu for near; with, then the Urdu for to be able to.
\end{quote}

\subsection*{You'll want to know: \ur{پانا}}
\textbf{\ur{پانا}} — \emph{pānā} — \textquotedblleft{}to manage to, to get\textquotedblright{}.

\begin{cousinweb}
\textbf{\ur{پانا}} (\emph{pānā}) is a tadbhava, a word inherited from Sanskrit and reshaped by sound change. It comes from \emph{prāpayati}, \textquotedblleft{}causes to be obtained, obtains\textquotedblright{} (CDIAL 8943).
\end{cousinweb}

\begin{grammarlens}[title={What you've built}]
One more word for this chapter.
\end{grammarlens}

\subsection*{Your turn}
\begin{itemize}
\raggedright
  \item \emph{Say it:} \emph{pānā}
  \item \emph{Say it:} \emph{pānā}, once more
  \item \emph{Say it:} say \emph{pānā}
  \item \emph{Recall:} say the Urdu for near; with, then the Urdu for to be able to, then say \emph{pānā} again
\end{itemize}

\begin{tcolorbox}[breakable,colback=teal!4,colframe=teal!35!black,title={Before you move on}]
What does \ur{پانا} mean? (\textquotedblleft{}To manage to, to get\textquotedblright{}.) Say it once more.
\end{tcolorbox}
